\clearpage
\section{Simplicity-to-mathematics correspondence}
\label{app:code-math}

This appendix maps the native audit verifier in the reference implementation to Section~\ref{sec:construction}.
The source is \href{https://github.com/BlockstreamResearch/damp/blob/d48f1833c5a5663ed2de48dfb2f7d762a9f3d19c/simf/lib/audit.simf}{\texttt{simf/lib/audit.simf}} at revision \texttt{d48f183}.
The numbered listings include that module directly.
The tables explain each source statement and group multiline statements together.
Blank lines have no semantic role; the tables group declaration delimiters with their enclosing statement.
The \texttt{.simf} files are SimplicityHL source compiled into the Simplicity programs executed by the covenant.
This is a source-level correspondence, not a machine-checked equivalence proof.

\lstdefinestyle{auditcode}{
  basicstyle=\ttfamily\fontsize{8.5}{10}\selectfont,
  numbers=left, numberstyle=\tiny, numbersep=8pt,
  xleftmargin=1.8em, frame=tb, columns=fullflexible,
  breaklines=true, breakatwhitespace=false, keepspaces=true,
  showstringspaces=false, tabsize=4, aboveskip=9pt, belowskip=9pt
}
\newcommand{\auditlisting}[2]{%
  \lstinputlisting[style=auditcode,firstline=#1,lastline=#2,firstnumber=#1]{../../simf/lib/audit.simf}%
}
\newcommand{\auditmap}[2]{\texttt{#1} & #2\\}
\newenvironment{auditmapping}{%
  \begingroup\small
  \renewcommand{\arraystretch}{1.12}%
  \begin{longtable}{@{}>{\raggedright\arraybackslash}p{.12\linewidth}>{\raggedright\arraybackslash}p{.85\linewidth}@{}}
  \toprule
  Source lines & Mathematical meaning or check\\
  \midrule
  \endhead
  \bottomrule
  \endfoot
}{%
  \end{longtable}\endgroup
}

\subsection{Notation and point encoding}

Several source names differ from the paper's symbols.
The local variable \texttt{c} is the point $C$, while \texttt{e} is the challenge $c$.
The local \texttt{h} is $H_A$, not the transaction digest $h$.
The argument \texttt{d} is $D$, not the deployment salt $d$; the salt and audit epoch come from \texttt{AUDIT\_DEPLOYMENT} and \texttt{AUDIT\_EPOCH}.
The variable \texttt{aux} holds $R$, and \texttt{root} is a field element $\rho$ used only to authenticate the native commitment's sign.

Write a compressed curve point as $(\eta,x_C)$, where $\eta$ is the SEC~1 parity bit defined in Section~\ref{sec:statement-proof}.
\texttt{Ge} denotes an affine point and \texttt{Gej} a Jacobian point.
The pair \texttt{(point, 1)} embeds an affine point with Jacobian coordinate $Z=1$.
\texttt{gej\_equiv} compares represented group elements, not their raw Jacobian coordinates.
The curve's base-field prime is $p_F=2^{256}-2^{32}-977$, distinct from the scalar order $q$.

\auditlisting{1}{19}
\begin{auditmapping}
  \auditmap{1-3}{State the module's purpose and import checked counter addition/subtraction. These imports do not implement a group equation.}
  \auditmap{5-6}{Declare \texttt{point\_ge} and unpack a compressed point into its parity bit and $x$ coordinate.}
  \auditmap{7}{Require $0\leq x<p_F$ by equality with field normalization. Reject a noncanonical coordinate.}
  \auditmap{8}{Decompress the curve point selected by $(\eta,x)$; \texttt{unwrap} rejects a non-point. This establishes a finite point on secp256k1.}
  \auditmap{9}{End the point-decoding helper.}
  \auditmap{11-12}{Declare \texttt{hash\_point} and unpack the point to serialize into a SHA-256 state.}
  \auditmap{13}{Select the SEC~1 prefix from parity.}
  \auditmap{14}{An even $y$ coordinate uses byte \texttt{02}.}
  \auditmap{15}{An odd $y$ coordinate uses byte \texttt{03}.}
  \auditmap{16}{Complete prefix selection.}
  \auditmap{17}{Append the one-byte SEC~1 prefix to the transcript.}
  \auditmap{18}{Append the 32-byte $x$ coordinate, completing $\operatorname{SEC1}(Q)$ for the supplied point $Q$.}
  \auditmap{19}{End the serialization helper.}
\end{auditmapping}

\subsection{Read the native commitment and validate proof fields}
\label{app:native-sign}

\auditlisting{21}{54}
\begin{auditmapping}
  \auditmap{21-31}{Declare the verifier with output index $i$, SEC~1 parity $\eta$, sign witness $\rho$, points $D,T_C,T_D$, responses $z_v,z_b$, and digest $\sha(R)$. These are public proof or encoding data, not $(v,b)$.}
  \auditmap{32}{Read the actual asset and amount field of output $i$ from the transaction. Reject a missing output.}
  \auditmap{33}{Require an explicit asset identifier, not an asset commitment.}
  \auditmap{34}{Require that identifier to equal the compiled regulated asset $A$.}
  \auditmap{35}{Require a confidential value commitment and unpack its native sign bit $\epsilon$ and coordinate $x_C$. Reject an explicit amount.}
  \auditmap{36}{Decode a candidate $C=(x_C,y_C)$ using the supplied SEC~1 parity $\eta$.}
  \auditmap{37}{Read its $y_C$ coordinate for the native-sign check.}
  \auditmap{39-40}{Document why the native Elements sign bit must not be interpreted as SEC~1 parity.}
  \auditmap{41}{Require the witness $\rho$ to be a canonical field element.}
  \auditmap{42}{Select a quadratic-residue target from the native sign bit.}
  \auditmap{43}{For prefix \texttt{08}, require a square root of $y_C$.}
  \auditmap{44}{For prefix \texttt{09}, require a square root of $-y_C$.}
  \auditmap{45}{Complete the target selection.}
  \auditmap{46}{Check $\rho^2=(-1)^\epsilon y_C\pmod {p_F}$. This ties the decoded point to the actual native encoding of $C$.}
  \auditmap{47}{Derive $H_A$ with the Elements asset-generator jet \texttt{hash\_to\_curve}.}
  \auditmap{48}{Read the compiled audit public key $P$.}
  \auditmap{49}{Decode $P$ as a canonical finite curve point.}
  \auditmap{50}{Decode the proof's handle $D$.}
  \auditmap{51}{Decode the proof commitment $T_C$.}
  \auditmap{52}{Decode the proof commitment $T_D$.}
  \auditmap{53}{Require $0\leq z_v<q$. Reject a noncanonical scalar encoding.}
  \auditmap{54}{Apply the same canonical-scalar check to $z_b$.}
\end{auditmapping}

The native prefixes \texttt{08} and \texttt{09} distinguish the quadratic character of $y_C$.
They do not mean even and odd $y_C$.
For secp256k1, $p_F\equiv3\pmod4$, so $-1$ is a quadratic nonresidue.
For nonzero $y_C$, exactly one of $y_C$ and $-y_C$ has a square root.
A finite point in this odd-prime-order group cannot have $y_C=0$.
Thus line 46 accepts only the sign selected by the native prefix, even though line 36 starts from a separately supplied SEC~1 parity bit.
Changing $C$ to $-C$ while retaining the same native commitment bytes fails this sign check.
The prover derives the witness $\rho$ from public curve coordinates; it does not reveal the opening of $C$.

\subsection{Construct the Fiat--Shamir challenge}

Let $t=\sha(\mathtt{DAMP/audit/autonomous/v3})$.
The challenge uses $\sha(t\mathbin\Vert t\mathbin\Vert m)$ with
\begin{align*}
 m={}&d\mathbin\Vert\operatorname{BE}_4(2)\mathbin\Vert\operatorname{BE}_8(e)
       \mathbin\Vert\operatorname{SEC1}(P)\mathbin\Vert h\mathbin\Vert\operatorname{BE}_4(i)\\
    &\mathbin\Vert\operatorname{asset}(A)\mathbin\Vert\operatorname{native}(C)
       \mathbin\Vert S\mathbin\Vert\mathtt{00}\mathbin\Vert\sha(R)\\
    &\mathbin\Vert\operatorname{SEC1}(D)\mathbin\Vert\operatorname{SEC1}(T_C)
       \mathbin\Vert\operatorname{SEC1}(T_D).
\end{align*}
Here $\operatorname{asset}(A)$ is the 32-byte consensus-order asset identifier, not its display-order hexadecimal text.
The encoding $\operatorname{native}(C)$ consists of the native prefix \texttt{08} or \texttt{09} and $x_C$, for 33 bytes in total.
The fixed role byte \texttt{00} is part of the transcript even though it is not a variable in $x$.
The \texttt{v3} tag names the challenge format; it does not change the domain word $g=2$.

\auditlisting{56}{84}
\begin{auditmapping}
  \auditmap{56-57}{Document the ordered domain and transaction fields.}
  \auditmap{58}{Start a SHA-256 context.}
  \auditmap{59-61}{Append $t$, whose hexadecimal value is the source constant.}
  \auditmap{62-64}{Append $t$ again, implementing the tagged-hash prefix $t\mathbin\Vert t$.}
  \auditmap{65}{Append the 32-byte deployment salt $d$.}
  \auditmap{66}{Append $\operatorname{BE}_4(g)$ with $g=2$.}
  \auditmap{67}{Append the 8-byte audit epoch $e$.}
  \auditmap{68}{Append $\operatorname{SEC1}(P)$.}
  \auditmap{69}{Append the introspected transaction digest $h=\texttt{sig\_all\_hash()}$.}
  \auditmap{70}{Append the 4-byte output index $i$.}
  \auditmap{71}{Append the 32-byte consensus asset identifier $A$.}
  \auditmap{72}{Select the native commitment prefix using the bit read from the transaction.}
  \auditmap{73}{Use \texttt{08} when $y_C$ is a quadratic residue.}
  \auditmap{74}{Use \texttt{09} when $y_C$ is a quadratic nonresidue.}
  \auditmap{75}{Complete the native-prefix selection.}
  \auditmap{76}{Append that one-byte native prefix.}
  \auditmap{77}{Append $x_C$, completing the transaction's exact 33-byte native commitment.}
  \auditmap{78}{Append $S=\sha(\mathrm{scriptPubKey}_i)$ from transaction introspection.}
  \auditmap{79}{Append the fixed output-role byte \texttt{00}.}
  \auditmap{80}{Append $\sha(R)$, calculated from the record bytes in lines 149-150 below.}
  \auditmap{81}{Append $\operatorname{SEC1}(D)$.}
  \auditmap{82}{Append $\operatorname{SEC1}(T_C)$.}
  \auditmap{83}{Append $\operatorname{SEC1}(T_D)$.}
  \auditmap{84}{Finalize SHA-256 and reduce its big-endian integer modulo $q$ to obtain challenge $c$. The source calls this scalar \texttt{e}.}
\end{auditmapping}

\subsection{Check the two group equations}

\noindent\begin{minipage}{\linewidth}
\auditlisting{86}{96}
\end{minipage}
\begin{auditmapping}
  \auditmap{86-87}{State the common-blinder relation and its two verification equations in source notation.}
  \auditmap{88}{Start the group-equality assertion for Equation~\ref{eq:verify-commitment}.}
  \auditmap{89}{\texttt{linear\_combination\_1} computes $z_vH_A+z_bG$. The generator $G$ is implicit in this jet.}
  \auditmap{90}{\texttt{scale} computes $cC$ and \texttt{gej\_ge\_add} adds $T_C$, giving $T_C+cC$.}
  \auditmap{91}{Require the two group elements to be equal, completing Equation~\ref{eq:verify-commitment}.}
  \auditmap{92}{Start the assertion for Equation~\ref{eq:verify-handle}.}
  \auditmap{93}{Compute $z_bP$, reusing the same response $z_b$ as line 89.}
  \auditmap{94}{Compute $cD$ and add $T_D$, giving $T_D+cD$.}
  \auditmap{95}{Require equality, completing Equation~\ref{eq:verify-handle}.}
  \auditmap{96}{Return only after both assertions succeed.}
\end{auditmapping}

The substitution proof in Section~\ref{sec:statement-proof} explains why honest responses pass these exact operations.
Neither assertion decrypts $R$, verifies a native range proof, or proves that this transaction was confirmed.

\subsection{Serialize recovery bytes and bind the native range proof}

\auditlisting{98}{130}
\begin{auditmapping}
  \auditmap{98}{Represent $R$ as $64+32+4+2=102$ bytes. This tuple splits raw bytes for SHA-256 jets; its boundaries are not the encryption-envelope field boundaries.}
  \auditmap{99}{Represent the range-proof body as $79\cdot64+4=5060$ bytes. The ten-byte header is fixed separately.}
  \auditmap{100}{Define one record as $(i,\eta,\rho,D,T_C,T_D,z_v,z_b,R,\mathrm{rangeBody})$. It contains no plaintext witness $(v,b)$.}
  \auditmap{101}{Allocate ten optional record slots, the output bound used by the verifier.}
  \auditmap{102}{Define fold state as the accepted count, previous index, finished flag, and recovery-manifest hash context.}
  \auditmap{104-105}{Declare \texttt{add\_aux} and unpack the four byte segments of $R$.}
  \auditmap{106}{Append bytes 0 through 63 of $R$.}
  \auditmap{107}{Append bytes 64 through 95 of $R$.}
  \auditmap{108}{Append bytes 96 through 99 of $R$.}
  \auditmap{109}{Append bytes 100 and 101 of $R$. Integers retain their fixed-width big-endian encoding.}
  \auditmap{110}{Return the hash context after all 102 bytes, without decoding the envelope.}
  \auditmap{112-114}{Define \texttt{range\_chunk}, which appends one 64-byte range-proof body chunk.}
  \auditmap{116-119}{Declare the range-binding check. Elements consensus verifies the range proof; this code authenticates the required profile.}
  \auditmap{120}{Unpack the 79 chunks and four-byte tail.}
  \auditmap{121}{Start a fresh SHA-256 context for the complete range proof.}
  \auditmap{122}{Append header bytes \texttt{60 3e}, selecting exponent zero and 63 mantissa bits.}
  \auditmap{123}{Append the eight-byte minimum value 1. Together the fixed header specifies $1\leq v\leq2^{63}$.}
  \auditmap{124}{Append all 79 body chunks in order.}
  \auditmap{125}{Append the final four bytes, completing the 5070-byte proof serialization.}
  \auditmap{126}{Begin a digest-equality assertion.}
  \auditmap{127}{Finalize the hash of the header and supplied body.}
  \auditmap{128}{Read the actual output range-proof hash by introspection.}
  \auditmap{129}{Require the hashes to agree, so a different body or profile cannot be substituted without a SHA-256 collision.}
  \auditmap{130}{End the range-binding helper. The Sigma equations do not replace consensus range-proof verification.}
\end{auditmapping}

The auxiliary tuple's 64-byte first segment crosses the envelope's ciphertext boundary.
The contract treats all segments as bytes.
It commits even unsupported or invalid recovery envelopes without checking whether they can be decrypted.

\subsection{Require one record per output and commit the ordered records}

For the accepted ordered records $(i_1,R_1),\ldots,(i_n,R_n)$, define
\[
 M=\sha\!\left(
   \operatorname{BE}_4(i_1)\mathbin\Vert R_1\mathbin\Vert\cdots
   \mathbin\Vert\operatorname{BE}_4(i_n)\mathbin\Vert R_n
   \mathbin\Vert\operatorname{BE}_4(n)\right).
\]
The final transaction output must have script bytes $\mathtt{6a20}\mathbin\Vert M$.
These are \texttt{OP\_RETURN}, a 32-byte push, and $M$.
This transaction-level commitment is distinct from the per-proof digest $\sha(R_j)$.

\auditlisting{132}{176}
\begin{auditmapping}
  \auditmap{132-133}{Declare the per-slot fold and unpack $(n,\mathrm{previous},\mathrm{finished},\mathrm{manifest})$.}
  \auditmap{134}{Select an empty or present record slot.}
  \auditmap{135}{An empty slot marks the list finished without changing its count or hash.}
  \auditmap{136-137}{Enter a present record and document the ordering requirement.}
  \auditmap{138}{Check whether an empty slot has already occurred.}
  \auditmap{139}{Reject a present record after a gap.}
  \auditmap{140}{Allow a record while the list remains contiguous.}
  \auditmap{141}{Complete the no-gap assertion.}
  \auditmap{142}{Unpack $(i,\eta,\rho,D,T_C,T_D,z_v,z_b,R,\mathrm{rangeBody})$.}
  \auditmap{143}{Authenticate the native range-proof profile of output $i$.}
  \auditmap{144}{Require $i>0$, excluding the verifier anchor at output zero.}
  \auditmap{145}{Inspect the previously accepted output index.}
  \auditmap{146}{The first present record has no predecessor constraint.}
  \auditmap{147}{Otherwise require $i_{j-1}<i_j$, rejecting duplicates and reordered records.}
  \auditmap{148}{Complete the predecessor check.}
  \auditmap{149-150}{Compute $\sha(R)$ over exactly the 102 witness bytes. This is the digest passed into Equation~\ref{eq:statement}.}
  \auditmap{151}{Reconstruct the statement and verify Equations~\ref{eq:verify-commitment} and \ref{eq:verify-handle} for this output.}
  \auditmap{152}{Append the four-byte index $i_j$ to the manifest context.}
  \auditmap{153}{Append $R_j$ to that same context.}
  \auditmap{154}{Increase $n$ by one with checked addition and retain $i_j$ as the predecessor.}
  \auditmap{155-157}{Return the updated state and close the branch and per-slot helper.}
  \auditmap{159}{Declare the complete record check with the expected regulated-output count.}
  \auditmap{160-161}{Fold all ten slots from count zero, no predecessor, an unfinished list, and an empty hash context.}
  \auditmap{162}{Require $n>0$.}
  \auditmap{163}{Require the number of accepted records to match the regulated-output summary. Each accepted record names a distinct actual regulated commitment.}
  \auditmap{165-166}{Document the transaction-level recovery commitment and the offchain decryption boundary.}
  \auditmap{167}{Append the four-byte count $n$ after all ordered $(i_j,R_j)$ pairs.}
  \auditmap{168}{Finalize $M$.}
  \auditmap{169}{Start hashing the required script with bytes \texttt{6a20}.}
  \auditmap{170}{Append $M$, so the expected script hash is $\sha(\mathtt{6a20}\mathbin\Vert M)$.}
  \auditmap{171}{Select the last output using checked subtraction of one from the output count.}
  \auditmap{172}{Begin the final script-hash equality check.}
  \auditmap{173}{Read the actual last output's script hash.}
  \auditmap{174}{Finalize the expected script hash.}
  \auditmap{175}{Require equality, binding the ordered recovery bytes and their count to the transaction.}
  \auditmap{176}{Return after all record and manifest checks succeed.}
\end{auditmapping}

The coverage argument also uses the caller's output scan.
In \texttt{simf/lib/assets.simf}, lines 45-60, each confidential regulated output contributes one to that scan; an explicit output contributes its positive explicit value.
In \texttt{simf/verifier/core.simf}, lines 77-108, the scan requires a positive contribution and an owner-bound script for every regulated output.
Let $J_c$ and $J_e$ be the confidential and explicit regulated-output sets.
The summary is $|J_c|+\sum_{i\in J_e}v_i$.
The record checks enforce distinct indices in $J_c$, so $n\leq|J_c|$.
Equality at line 163 therefore forces $J_e$ to be empty and every member of $J_c$ to have exactly one accepted record.
The variable names containing \texttt{amount} in the shared scan must not be read as recovered confidential balances.

\subsection{Reachability from the spending program}

The depth-four entrypoint below calls the shared audit module after checking the anchor, input policy proofs, and output scripts.
The depth-five and depth-six entrypoints have the same audit call at line 24, each using its own policy-proof verifier.
These line numbers refer to \texttt{simf/verifier\_d4.simf}, not the shared audit module.

\lstinputlisting[style=auditcode,firstline=1,lastline=25,firstnumber=1]{../../simf/verifier_d4.simf}
\begin{auditmapping}
  \auditmap{1-2}{Document the entrypoint and import the shared \texttt{verify\_records} function.}
  \auditmap{3-11}{Import transaction-shape, counting, and budget helpers.}
  \auditmap{12}{Select the depth-four policy-proof checker.}
  \auditmap{14}{Enter the SimplicityHL spending program.}
  \auditmap{15}{Consume the fixed 4096-word witness padding, requiring zero words. This is an execution-budget device, not a Sigma equation.}
  \auditmap{16}{Check input/output zero as the verifier anchor and bind the initiating holder input to its owner.}
  \auditmap{17}{Scan regulated inputs and count those that need policy proofs; do not decrypt their amounts.}
  \auditmap{18-21}{Check the supplied input policy proofs with a bounded ten-slot fold.}
  \auditmap{22}{Require every regulated input to have been checked.}
  \auditmap{23}{Check regulated output scripts and obtain the summary used by the coverage argument above.}
  \auditmap{24}{Pass \texttt{witness::AUDIT\_RECORDS} to the shared verifier. This reaches the range, statement, two-equation, and manifest checks.}
  \auditmap{25}{End the entrypoint only after all checks succeed.}
\end{auditmapping}

\subsection{Operations outside Simplicity}

\begin{samepage}
The prover computes $C=vH_A+bG$, $D=bP$, $T_C,T_D$ and $z_v,z_b$ offchain.
The covenant checks their equations onchain at the lines identified above; the prover does not supply the secret scalars.
The Rust module \texttt{damp-core::native\_audit::proof} in the reference implementation constructs and rechecks these proof messages \cite{pgc-implementation}.
\end{samepage}

The recovery module's \texttt{AuditOutput::seal} constructs $R$ using ephemeral ECDH and authenticated encryption.
\texttt{VerifiedAuditProof::open\_recovery} decrypts it with the issuer key, then checks the recovered opening against both $C$ and $D$.
Neither operation appears in the Simplicity listing.
Likewise, \texttt{recover\_bounded} computes Equation~\ref{eq:recovery-target} and searches offchain:
\[
 C-s^{-1}D
 =vH_A+bG-s^{-1}(bsG)
 =vH_A.
\]
The inverse $s^{-1}$ is a scalar inverse modulo $q$ and exists for a nonzero audit secret.
The equation removes the blinder but does not solve the discrete logarithm for $v$.
Elements consensus verifies native range proofs, and the chain-validation/reporting layer checks chain inclusion.
The appendix does not attribute either guarantee to the two Sigma assertions.
